\section{Prolongación monodrómica y cálculo a profundidad prescrita}
\label{sec:generacion-monodromica-certificada-rev7}
\index{monodromía nonádica}\index{generación estructural}

La construcción parte del estado enriquecido producido por
\(\mathrm{APP}\to\mathrm{TRIT}\to\mathrm{TPK}\): catálogo activo de
\(468\) emisiones, proyección ternaria, acción \(D_3\), calibre de
orientación, registro de transiciones de rango pleno, caracteres de
clausura, propagación y autoescala, cilindros ternarios, fase nonádica y
sombra Paley--Witt. El registro reconstruye de manera única los
levantamientos \(L_0,L_1\) antes de abrir la carta arquimediana. Ninguna
cifra de \(\pi\), \(e\) o \(\varphi\) se recibe como entrada del generador.

\subsection{Primera prolongación y frontera coinductiva}

El censo exhaustivo selecciona las palabras
\[
 w_6^\pi=010211,\qquad
 w_6^e=201101,\qquad
 w_6^\varphi=121200.
\]
Con vectores fila sobre \(\mathbb F_3^6\), las dos elevaciones reconstruidas
son
\[
L_0=
\begin{pmatrix}
2&2&2&1&2&1\\
2&1&2&2&1&1\\
1&0&0&1&0&2\\
0&0&1&1&1&0\\
2&2&2&0&2&1\\
2&1&2&1&0&0
\end{pmatrix},
\qquad
L_1=
\begin{pmatrix}
2&1&2&0&2&2\\
1&2&1&1&2&0\\
1&2&1&1&1&2\\
0&1&0&0&2&1\\
2&0&2&1&0&1\\
0&2&2&2&1&1
\end{pmatrix}.
\]
Sean \(B^{\mathsf C}\), \(B^{\mathsf P}\) y \(B^{\mathsf A}\) las
biografías estructurales de clausura, propagación y autoescala que se
publican a continuación, descompuestas en bloques
\(b_0^\chi|\cdots|b_4^\chi\). Formemos
\[
 \begin{aligned}
 X_0&=(b_0^{\mathsf C},b_1^{\mathsf C},
       b_0^{\mathsf P},b_1^{\mathsf P},
       b_0^{\mathsf A},b_1^{\mathsf A})^T,\\
 Y_0&=(b_1^{\mathsf C},b_2^{\mathsf C},
       b_1^{\mathsf P},b_2^{\mathsf P},
       b_1^{\mathsf A},b_2^{\mathsf A})^T,
 \end{aligned}
\]
y, análogamente, \(X_1,Y_1\) con las transiciones
\(b_2\to b_3\) y \(b_3\to b_4\). Cada bloque se obtiene por la proyección
ternaria del estado producido tras la actualización completa
\[
 x_{j+1}^{\chi}
 =\mathcal U_{9j+9}\circ\cdots\circ\mathcal U_{9j+1}(x_j^{\chi}),
 \qquad
 b_j^{\chi}=\Pi_6(x_j^{\chi}),
 \qquad
 \mathcal U_t=\mathsf{Upd}_t\circ\mathsf{Tra}_t\circ\mathsf{Sel}_t.
\]
Así, el registro que precede a los levantamientos queda publicado
íntegramente por las cuatro matrices
\begingroup\scriptsize
\begin{equation}
\begin{aligned}
X_0&=\begin{pmatrix}
0&1&0&2&1&1\\0&1&2&2&2&2\\2&0&1&1&0&1\\
1&2&1&2&2&1\\1&2&1&2&0&0\\1&1&2&2&0&2
\end{pmatrix},&
Y_0&=\begin{pmatrix}
0&1&2&2&2&2\\0&1&0&2&1&1\\1&2&1&2&2&1\\
1&0&2&0&1&1\\1&1&2&2&0&2\\1&2&1&0&2&0
\end{pmatrix},\\[1ex]
X_1&=\begin{pmatrix}
0&1&0&2&1&1\\0&0&2&1&1&1\\1&0&2&0&1&1\\
0&1&2&2&2&2\\1&2&1&0&2&0\\0&1&0&2&1&0
\end{pmatrix},&
Y_1&=\begin{pmatrix}
0&0&2&1&1&1\\1&1&0&2&2&1\\0&1&2&2&2&2\\
1&0&2&0&1&1\\0&1&0&2&1&0\\0&1&0&2&0&0
\end{pmatrix}.
\end{aligned}
\label{eq:registro-transiciones-tpk-explicito-rev10}
\end{equation}
\endgroup
Sus filas son salidas sucesivas del endomorfismo TPK sobre las tres semillas
seleccionadas por los invariantes de fibra; no son coeficientes impuestos a
\(L_0,L_1\) ni datos tomados de las expansiones arquimedianas.

\begin{theorem}[Reconstrucción única de los levantamientos y del calendario]
\label{thm:reconstruccion-lifts-calendario-rev10}
En \(\mathbb F_3\) se cumple
\[
 \det X_0=1,
 \qquad \det X_1=-1,
 \qquad \operatorname{rank}(I_6\otimes X_i)=36.
\]
Por consiguiente, los sistemas
\[
 X_0L_0=Y_0,
 \qquad X_1L_1=Y_1
\]
poseen las soluciones únicas \(L_i=X_i^{-1}Y_i\), que son exactamente las
matrices impresas arriba. Entre los \(16\) calendarios binarios de longitud
\(4\), \(0011\) es el único que conserva simultáneamente las tres
biografías. Las \(144\) sustituciones unitarias posibles en ambas matrices
destruyen esa compatibilidad triple.
\end{theorem}

\begin{proof}
Los determinantes y rangos se calculan por eliminación exacta módulo \(3\).
La invertibilidad de \(X_i\) da existencia y unicidad de \(L_i\); los
\(16\) calendarios y las \(144=2\cdot6\cdot6\cdot2\) perturbaciones forman
dos censos finitos exhaustivos. Todas sus entradas son palabras, transiciones
y residuos del estado enriquecido anterior a la evaluación real. Los nombres
convencionales de las constantes no intervienen en ninguno de esos cálculos.
\end{proof}

La secuencia
\[
u_1=u_0L_0,\qquad u_2=u_1L_0,\qquad
u_3=u_2L_1,\qquad u_4=u_3L_1
\]
produce
\[
\begin{aligned}
w_{30}^{\pi}&=010211|012222|010211|002111|110221,\\
w_{30}^{e}&=201101|121221|102011|012222|102011,\\
w_{30}^{\varphi}&=121200|112202|121020|010210|010200.
\end{aligned}
\]
La matriz
\[
R_{36}=
\begin{pmatrix}
222220\\021222\\102011
\end{pmatrix}
\]
es la primera frontera coinductiva. Las matrices \(L_0,L_1\) son resultados
de la reconstrucción interna del
\cref{thm:reconstruccion-lifts-calendario-rev10} y se aplican antes de abrir
los testigos decimales. El dominio del teorema es la única cadena generadora
\(\mathrm{APP}\to\mathrm{TRIT}\to\mathrm{TPK}\to\) estado enriquecido.
APP aporta sus dos hojas aritméticas; TRIT orienta el régimen; TPK construye
y conserva los campos de ruta, carry, frontera y memoria sobre los que actúan
los levantamientos. Pedir que la proyección visible de APP contenga por sí
sola el estado terminal confunde dos dominios consecutivos de esa misma
cadena y no convierte los levantamientos en datos introducidos a mano.

\subsection{Selección finita y construcción de los caracteres}

Sean \(\mathcal C\) el multiconjunto de semillas, \(\mathcal U\) el catálogo
de emisiones y \(W=\Psi(\mathcal U)\subset\mathbb F_3^6\), con
\[
 \mathcal C\xrightarrow{T_{54}^{\rm seed}}\mathcal U
 \xrightarrow{\ \Psi\ }W\lhook\joinrel\longrightarrow\mathbb F_3^6,
 \qquad
 \Psi(U_6)=((-u_j)\bmod3)_{j=1}^6.
\]
La enumeración exacta da
\[
 \sum_{U\in\mathcal U}\operatorname{mult}(U)=104\,976,\qquad
 |\mathcal U|=468,\qquad |W|=243,\qquad|\mathbb F_3^6|=729.
\]
Estos cuatro números cuentan objetos tipados distintos.

La acción \(D_3\) sobre \(W\) produce \(43\) órbitas. La órbita de clausura
es la única cuyas seis orientaciones poseen cinco elevaciones, masa \(1008\)
y multiplicidades \(432+4\cdot144\). Entre las órbitas radicales unitarias,
el censo orientado de la clausura determina de modo único la propagación y el
censo \((2,2,2)\) determina de modo único la autoescala. El calibre
\(\operatorname{diag}_c=6\), con orientación ES, fija una palabra en cada
órbita. En consecuencia, la terna
\[
 (w_6^\pi,w_6^e,w_6^\varphi)
 =(010211,201101,121200)
\]
se obtiene antes de construir cualquier intervalo arquimediano.

Los tres caracteres también se obtienen de la estructura seleccionada. Para
la clausura, la matriz de incidencia
\[
 A_W=\begin{pmatrix}
 0&1&1&1&1&1\\
 1&0&1&2&2&1\\
 1&1&0&1&2&2\\
 1&2&1&0&1&2\\
 1&2&2&1&0&1\\
 1&1&2&2&1&0
 \end{pmatrix}
 \quad\text{satisface}\quad A_W^2=2I_6=-I_6
\]
y fija la inversión orientada. Sobre la pentafibra, el simetrizador que
olvida la etiqueta de la elevación es
\[
 P_5=\frac15\mathbf1\mathbf1^T.
\]
La conservación \(\mathbf1^T\lambda=1\) y la invariancia
\(P_5\lambda=\lambda\) fuerzan \(\lambda=(1/5)\mathbf1\), luego
\(t_1=1/5\). Las cuatro
composiciones por
\[
 t\star s=\frac{t+s}{1-ts}
\]
dan \(t_4=120/119\). La ecuación de retorno a la diagonal,
\[
 \frac{120/119-1/q}{1+(120/119)/q}=1,
\]
posee la única solución entera positiva \(q=239\). La identidad gaussiana
\begin{equation}
 \label{eq:generacion-monodromica-identidad-gaussiana-rev7}
 (5+i)^{16}(239-i)^4=-681386607803576157184
\end{equation}
reconoce después la media vuelta. Así, la expresión
\[
 16\arctan(1/5)-4\arctan(1/239)
\]
es la evaluación exacta del carácter de clausura, no su selector.

Para la propagación, la composición de caminos y la unidad conservada dan
\[
 P_{s+t}=P_sP_t,\qquad P_0=1,\qquad [t]P_t=1.
\]
Si \(P_t=\sum_{n\geq0}a_nt^n\), la comparación del coeficiente de \(s^nt\)
fuerza
\[
 (n+1)a_{n+1}=a_n,\qquad a_0=1,
\]
y, por inducción, \(a_n=1/n!\). La igualdad \(P'_t=P_t\) y el valor
\(P_1=e\) son reconocimientos analíticos posteriores de esta ley formal.

Para la autoescala, el calendario \((+,-,0)\) induce
\[
 F_{\rm av}=\begin{pmatrix}0&1\\1&1\end{pmatrix}.
\]
Su único rayo positivo invariante \(v=(1,x)^T\) satisface
\[
 F_{\rm av}v=xv,\qquad x^2-x-1=0,\qquad x>0.
\]
Los cocientes consecutivos de Fibonacci proporcionan horquillas racionales
encajadas para \(x\); la ecuación \(x=1+1/x\) es la coordenada escalar de
esta ecuación propia, no una condición impuesta desde fuera.

\begin{proposition}[Separación entre selección, evaluación y reconocimiento]
\label{prop:generacion-monodromica-separacion-causal-rev7}
Para cada canal \(\chi\), la composición causal es
\[
 \widehat X
 \xrightarrow{\operatorname{Sel}_{\rm TPK}}
 X_\chi
 \xrightarrow{\operatorname{ev}_{\mathbb R}}
 \mathbb R
 \xrightarrow{\operatorname{Rec}}
 \{\pi,\varphi,e,\alpha\}.
\]
\(\operatorname{Sel}_{\rm TPK}\) utiliza el censo, la órbita, el calibre,
los levantamientos, el acarreo y el carácter interno. La evaluación contrae
los cilindros compatibles. El reconocimiento asigna el nombre convencional
al valor ya construido. Ningún prefijo objetivo interviene en las dos primeras
flechas.
\end{proposition}

\begin{proof}
La unicidad de las tres órbitas se decide mediante predicados enteros del
catálogo. Los parámetros \(5,4,239\), la recurrencia
\((n+1)a_{n+1}=a_n\) y la matriz \(F_{\rm av}\) se deducen de esas órbitas y
de las leyes de unidad, composición e incidencia. La coordenada de
\(\alpha\) procede después del triple interno mediante el registro firmado
\(U_{12}\), el sello \(K\) y el acarreo dodecafásico. Las cotas racionales se
calculan únicamente para evaluar y certificar la sección ya cerrada. La
comparación decimal se ejecuta después; por tanto, no puede actuar como
entrada causal.
\end{proof}

\subsection{Extensión de supervivientes y certificación posterior}

Sea \(\mathcal H_m^\chi\) el conjunto de historias enriquecidas admisibles
del canal \(\chi\) a profundidad \(m\), y sea
\(p_{n,m}:\mathcal H_n^\chi\to\mathcal H_m^\chi\) el truncamiento. Se define
\[
 \operatorname{Surv}_{m}^{(N),\chi}
 =p_{N,m}(\mathcal H_N^\chi),\qquad
 \operatorname{Surv}_{m}^{\chi}
 =\bigcap_{N\ge m}\operatorname{Surv}_{m}^{(N),\chi},
\]
y la sección numérica por
\[
 x_\chi\in X_\chi:=
 \varprojlim_m\operatorname{Surv}_m^\chi.
\]
El lector y el estado no se identifican: una presentación
\((x_\chi,L_\chi)\) publica el valor, mientras el estado conserva también
orientación, acarreo, frontera, hoja, incidencia excepcional y genealogía.
En la frontera R36, \(\chi_{\rm HMT}\) y el residuo interno
\(\varepsilon_K^\chi\) son ya coordenadas del estado. La división APP y el
selector de fase producen
\[
 d_K^\chi=\operatorname{Dig}_{729}(\varepsilon_K^\chi),\qquad
 \varepsilon_{K+1}^\chi=729\varepsilon_K^\chi-d_K^\chi,qquad
 N_{K+1}^{\operatorname{Ext}}=729N_K+d_K^\chi,
\]
y la selección, el transporte y la actualización
\[
 \operatorname{Sel}_K^\chi(x)
 :=\operatorname{Sel}\!\left(
 x,M_{\rm ph}(\rho_9(K+1))\right),
 \qquad
 \operatorname{Tra}_K^\chi(y,d):=\operatorname{Lift}(T_d)(y),
 \qquad
 \operatorname{Upd}_K^\chi:=\operatorname{Rec}_K^\chi\circ
 \operatorname{Mem}_K^\chi,
\]
\[
 \mathcal U_K^\chi(x)
 :=\operatorname{Upd}_K^\chi\!\left(
  \operatorname{Tra}_K^\chi\!\left(
   \operatorname{Sel}_K^\chi(x),d_K^\chi
  \right)\right)
\]
depende sólo del estado anterior. La relación completa
\(\operatorname{Ext}\) puede ser multivaluada; la subrelación
\(\operatorname{Ext}^{\chi,\rightarrow}\), anclada en el estado R36 que ya
porta el carácter, es el grafo de \(\mathcal U_K^\chi\). Por ello fija
una sección global compatible sin afirmar unitariedad de cada fibra local de
\(\operatorname{Ext}\). No es una intersección retrospectiva definida por
\(P_{\rm ar}\), y no lee cifras futuras ni horquillas.

Si \(N_K\) es el prefijo ternario y \(L_K=6K\), sus \(729\) subcilindros son
\[
 I_{K,u}=
 \left[
 \frac{729N_K+u}{3^{L_K+6}},
 \frac{729N_K+u+1}{3^{L_K+6}}
 \right),
 \qquad 0\le u<729.
\]
Para la horquilla racional exacta \([\ell_{\chi,K},h_{\chi,K}]\) construida
por la recurrencia interna del carácter \(\chi\),
\[
\begin{aligned}
j_{\min}
 &=\max\!\left(729N_K,
   \left\lfloor\ell_{\chi,K}3^{L_K+6}\right\rfloor\right),\\
j_{\max}
 &=\min\!\left(729N_K+728,
   \left\lceil h_{\chi,K}3^{L_K+6}\right\rceil-1\right).
\end{aligned}
\]
La comprobación arquimediana posterior es
\[
 j_{\min}=j_{\max}=N_{K+1}^{\operatorname{Ext}}.
\]
Certifica que la partición es exhaustiva y localiza el superviviente ya
producido por \(\operatorname{Ext}^{\chi,\rightarrow}\); no lo selecciona.
La horquilla no se lee de una expansión decimal: se calcula
desde la pentafibra y el retorno unitario, desde la recurrencia de propagación
o desde la incidencia \(F_{\rm av}\), según el canal. Por tanto, la
intersección es un certificado de publicación, no una regla de descendencia.

\begin{theorem}[Generación coinductiva correlacionada de profundidad arbitraria]
\label{thm:generacion-monodromica-conjunta-rev7}
En la estructura enriquecida TPK construida en los teoremas anteriores, los
tres caracteres derivados y el cierre dodecafásico construyen a profundidad
arbitraria la familia correlacionada de secciones
\(x_\pi,x_\varphi,x_e,x_\alpha\) que satisface:
\begin{enumerate}
 \item el barrido
 \[
 104\,976\longrightarrow468\longrightarrow243
 \]
 separa semillas, emisiones y palabras visibles;
 \item los levantamientos
 \[
 w_6\xrightarrow{L_0}w_{12}\xrightarrow{L_0}w_{18}
 \xrightarrow{L_1}w_{24}\xrightarrow{L_1}w_{30}
 \]
 preceden a la primera frontera coinductiva \(R_{36}\);
 \item en cada nivel,
 \(\operatorname{Ext}^{\chi,\rightarrow}\) produce primero el hijo
 compatible y su truncamiento coincide con el estado del nivel anterior; los
 \(729\) subcilindros numéricos forman una partición ordenada y la horquilla
 racional localiza después el hijo ya generado;
 \item la fase cumple \(g(K+9)=g(K)\), mientras el bloque, el índice de
 cilindro, el prefijo y la sombra de Witt cambian;
 \item los lectores publican, respectivamente, \(\pi\), \(\varphi\), \(e\)
 y \(\alpha\), y la regla de extensión carece de profundidad terminal.
\end{enumerate}
\end{theorem}

El alcance del enunciado es coinductivo y de profundidad arbitraria. Para los
cuatro canales, la unicidad de la sección forward y su compatibilidad bajo
truncamiento proceden de la actualización de estado. No se afirma que toda
fibra local de \(\operatorname{Ext}\) sea unitaria. Las corridas a cualquier
profundidad finita son sólo testigos de regresión y publicación.

\begin{proof}
La selección finita de
\cref{prop:generacion-monodromica-separacion-causal-rev7} conserva las cinco
regiones de la órbita de clausura antes de abrir la carta real. La pendiente
primitiva \(1/5\), la composición de las cuatro ramas acompañantes y la
condición de retorno unitario fuerzan, en ese orden, \(120/119\) y \(q=239\).
La identidad
\eqref{eq:generacion-monodromica-identidad-gaussiana-rev7} reconoce entonces
la primera media vuelta positiva, de modo que el evaluador posterior es
\[
 \pi_{\rm HMT}
 =16\arctan(1/5)-4\arctan(1/239).
\]
Las arctangentes se evalúan por series alternadas con cotas racionales.

En propagación, la ley \(P_{s+t}=P_sP_t\), la unidad \(P_0=1\) y la
normalización lineal fuerzan primero \((n+1)a_{n+1}=a_n\). Sólo después se
reconocen \(P'_t=P_t\) y
\[
 e_{\rm HMT}=\sum_{n\ge0}\frac1{n!}.
\]
En autoescala, el calendario produce primero
\[
 F_{\rm av}=\begin{pmatrix}0&1\\1&1\end{pmatrix}
\]
y su rayo de Perron. La ecuación \(x=1+1/x\) es su coordenada escalar
posterior.

La actualización que sólo lee el estado anterior produce en cada canal el
hijo compatible a partir del carácter y el residuo internos presentes. Por
inducción, dos secciones forward con el mismo estado R36 coinciden a toda
profundidad. En los tres canales con lector real desarrollado aquí, la
aritmética racional produce después intervalos de anchura arbitrariamente
pequeña: la partición exacta de los \(729\) subcilindros comprueba que el hijo
ya generado queda dentro de la celda decimal publicada. El cierre firmado
\(U_{12}\to K\) prolonga la cuarta coordenada y los cuadrados de truncamiento
de sus dos vías internas conmutan, de modo que
\(\alpha_A=\alpha_B=\alpha_{\rm HMT}\). Las ejecuciones de mil y diez mil
cifras son proyecciones finitas, no el criterio del teorema.
\end{proof}

La extensión a una sección del límite inverso se rige por el
\cref{thm:generacion-supervivencia-pi-e-phi-rev6}. El teorema anterior
descarga, para \(\pi,\varphi,e,\alpha\), la no vacuidad, la compatibilidad y
la unicidad global de cada sección forward exigidas por ese principio
abstracto, sin postular unitariedad puerta a puerta de
\(\operatorname{Ext}\). Esta formulación distingue la regla coinductiva de
profundidad arbitraria de sus evaluaciones finitas de mil y diez mil cifras,
sin alterar el retorno de fase ni el transporte de memoria.

\begin{corollary}[Profundidad arbitraria y objeto decimal infinito]
\label{cor:generacion-nonadica-profundidad-arbitraria-rev10}
Para todo \(N\in\mathbb N_{>0}\), la misma construcción publica un único
prefijo decimal canónico \(d_N^\chi\) de longitud \(N\), con
\(d_{N+1}^\chi|_N=d_N^\chi\), para
\(\chi\in\{\pi,\varphi,e,\alpha\}\). La familia compatible
\((d_N^\chi)_{N\ge1}\) determina la expansión infinita en el límite
inverso. Mil, diez mil o un millón de cifras son ejecuciones finitas del
mismo algoritmo parametrizado; ninguna de ellas limita el alcance del
teorema.
\end{corollary}

\begin{proof}
La actualización \(\operatorname{Upd}_\chi\) produce primero el cilindro
compatible. La horquilla racional del carácter tiene anchura arbitrariamente
pequeña y certifica después que ese cilindro intersecta un único miembro de
la partición publicada de \(729\) subcilindros. Dado \(N\), existe \(m\) con
\(3^{-6m}<10^{-N}\); el cilindro ya producido queda entonces contenido en una
celda decimal canónica de longitud \(N\). La compatibilidad de truncación
identifica esas celdas como las coordenadas finitas de una sección forward.
\end{proof}

\begin{theorem}[Generación coinductiva correlacionada y doble proyección]
\label{thm:estado-coinductivo-doble-proyeccion-rev10}
Sea \(\mathscr S_\infty\) el producto fibrado de las cuatro secciones
forward correlacionadas de profundidad arbitraria generado por
\[
 \mathrm{APP}\longrightarrow\mathrm{TRIT}\longrightarrow\mathrm{TPK}
 \longrightarrow R_{36}\longrightarrow G_9.
\]
Desde \(\mathscr S_\infty\) parten dos publicaciones tipadas:
\[
 \mathscr S_\infty\xrightarrow{\mathcal P_{\rm Arq}}\mathbb R^4,
 \qquad
 \mathscr S_\infty\xrightarrow{\mathcal P_{\rm Exc}}
 W_{12}\longrightarrow G_{24}\longrightarrow\Lambda_{24}
 \longrightarrow V^\natural
 \longrightarrow\operatorname{Aut}(V^\natural)=\mathbb M.
\]
La primera contrae cilindros compatibles y publica correlacionadamente
\((\pi_{\rm HMT},\varphi_{\rm HMT},e_{\rm HMT},\alpha_{\rm HMT})\); la
segunda conserva la incidencia de frontera y publica la cadena excepcional.
En consecuencia, \(\mathbb R\) es la sombra arquimediana de la primera
proyección y no una entrada que seleccione la genealogía. Machin,
Perron--Fibonacci, el semigrupo exponencial y las denominaciones
convencionales actúan después de la salida HMT.
\end{theorem}

\begin{proof}
El \cref{thm:generacion-monodromica-conjunta-rev7} y el
\cref{cor:generacion-nonadica-profundidad-arbitraria-rev10} construyen la
familia compatible antes de la evaluación. El cierre dodecafásico publica
\(\alpha_{\rm HMT}\) desde su triple R36 correlacionado y el registro
\(U_{12}\to K\); la vía
\(E_{21/22}\) conserva su dominio analítico propio. La proyección
arquimediana olvida incidencia al contraer los cilindros, mientras la
proyección excepcional la transporta por los entrelazadores Witt--Golay--
Leech hasta \(V^\natural\). Ambas proceden del mismo antecedente y ninguna
puede actuar retroactivamente como generador de éste.
\end{proof}

\subsection{Regiones de clausura}

Las cinco regiones no son cinco valores de \(\pi\), sino cinco genealogías
con una proyección visible común. Su multiplicidad es
\[
 432+4\cdot144=1008,
\]
y sus papeles son
\[
 3_{++}+1_{--}+1_{\perp}.
\]
La identidad regional se verifica por la cuádrupla
\[
 (U_6,E_{\rm sig},\text{multiplicidad},\text{papel}),
\]
no sólo por el censo \(3/1/1\). En particular,
\[
\begin{array}{c|c|c|c}
U_6&E_{\rm sig}&\text{multiplicidad}&\text{papel}\\ \hline
501|614|498|169|272|272&555555&432&\perp\\
810|923|870|169|272|272&339555&144&++\\
810|983|810|169|272|272&393555&144&++\\
870|923|810|169|272|272&933555&144&++\\
870|923|810|418|674|521&933717&144&--
\end{array}
\]
La tabla fija conjuntamente palabra decimal enriquecida, firma de emisión,
multiplicidad y papel. En particular, el ancla de multiplicidad (432) es
transversal y la firma (933717) es el retorno mutado; el censo (3/1/1)
por sí solo no determina esos nombres.

\subsection{Retorno de fase y cambio de memoria}

En la novena fase,
\[
 B_9=(100100\mid020112\mid010122)
\]
admite tres orientaciones que conservan el eje \(\varphi\) y el agregado
\(\pi+e\):
\[
\begin{aligned}
 &(101222\mid112221\mid112002),\\
 &(102221\mid111222\mid112002),\\
 &(111221\mid102222\mid112002).
\end{aligned}
\]
Las tres sobreviven una frontera y conservan
\[
 u^{\varphi}=112002,
 \qquad u^\pi+u^e=210110.
\]
La frontera siguiente selecciona
\[
 B_{10}=(101222\mid112221\mid112002).
\]
Por tanto,
\[
 \#\operatorname{Surv}_1(B_9)=3,
 \qquad
 \#\operatorname{Surv}_2(B_9)=1.
\]
El índice diez ocupa de nuevo la fase uno, pero no el estado anterior. Ésta
es la monodromía: retorno de fase con transformación de memoria.

En términos del estado completo
\[
 \widehat x_K=(B_K,C_K,p_K,c_K,\Sigma_K,W_K,g(K)),
\]
la vuelta es la holonomía única
\[
 \begin{gathered}
 \Gamma_9=\operatorname{Hol}_{C_9}(\nabla_{\rm TPK}),
 \qquad g(K+9)=g(K),\\
 q_{\rm mem}\Gamma_9=q_{\rm mem}+1,
 \qquad \Gamma_9\widehat x_K\ne\widehat x_K.
 \end{gathered}
\]
En realización isométrica, la misma continuidad satisface
\[
 U_{\Gamma_9}J=JT+\eta,\qquad
 I-T^*T=\eta^*\eta.
\]
El acarreo se compone por
\[
 \operatorname{Carr}(\gamma_2\gamma_1)
 =\operatorname{Carr}(\gamma_2)H(\gamma_1)
 +Y(\gamma_2)\operatorname{Carr}(\gamma_1),
\]
y la telescopía conserva el terminal finito:
\[
 S_0=\sum_{r=0}^{L-1}M_9^{*r}D_9^*D_9M_9^r
     +M_9^{*L}S_0M_9^L.
\]
Las ventanas \(27\), \(54\) y \(108\) son las concatenaciones
\[
 \Gamma_{27}=\Gamma_9^3,\qquad
 \Gamma_{54}=\Gamma_9^6,\qquad
 \Gamma_{108}=\Gamma_9^{12};
\]
transportan una memoria acumulada de \(3\), \(6\) y \(12\) vueltas,
respectivamente. No constituyen reinicios del algoritmo.

\subsection{Censo finito de transiciones y generación decimal}

La ejecución a profundidad mil produce, por canal,
\[
396\ \text{transiciones},\quad
2376\ \text{trits},\quad
43\ \text{vueltas completas},\quad
387\ \text{pares }(K,K+9),
\]
y \(1000\) cifras decimales. En cada transición, la actualización
\(\operatorname{Ext}^{\chi,\rightarrow}\) produce el superviviente y el
barrido exhaustivo posterior de los \(729\) subcilindros certifica su
publicación. Los testigos externos se consultan después de la generación y se
usan sólo para el reconocimiento.

La independencia del blanco se obtiene porque el mapa generador recibe
únicamente el catálogo finito de \(468\) emisiones, las matrices
\(L_0,L_1,A_W,F_{\rm av}\) y aritmética racional. Ninguna de sus funciones
de transición contiene prefijos de \(\pi\), \(e\), \(\varphi\) ni una
operación que evalúe esas constantes. Las horquillas se deducen de los
caracteres construidos y actúan únicamente en la evaluación posterior.

Tras las 396 etapas de prolongación y certificación, el cilindro final
\[
 \left[\frac{N}{3^{2376}},\frac{N+1}{3^{2376}}\right)
\]
queda contenido en una única celda decimal de profundidad \(10^{-1000}\).
Las mil cifras son, por ello, un testigo finito de la ley parametrizable, no
su profundidad máxima.

El argumento uniforme construye un cilindro para cualquier profundidad
finita prescrita. El censo de \(396\) etapas sólo comprueba la instancia de
mil cifras: conserva separadas la generación causal y la comparación decimal
posterior, y no reemplaza el cuantificador del
\cref{thm:generacion-monodromica-conjunta-rev7}.
